\section{Justificación}

% Cuál es el problema que se investiga

% Cual es la emergencia del problema de estudio

% Que tan viable es la realización de esta investigación

% A qué segmento de la población beneficia el estudio

% Qué beneficios metodológicos tiene el estudio

% Que beneficio personal, profesional o disciplinario se obtendrá al realizar el estudio

Las personas con discapacidad visual pueden enfrentar problemas de salud mental relacionados con la adaptación, la ansiedad o la depresión debido a su condición [3].  

\\

Se ha observado que la capacidad de adaptación en la movilidad independiente de las personas con discapacidad visual varía significativamente según el momento en que se presenta esta condición. 

\\

Aquellos que nacen con discapacidad visual, ya sea total o parcial, suelen desarrollar desde temprana edad estrategias sensoriales y cognitivas que les permiten desenvolverse con mayor autonomía en su entorno. Aun así, enfrentan limitaciones en su desarrollo y orientación espacial; especialmente al interactuar con entornos nuevos o no 
accesibles, donde, al presenciar alguna situación de emergencia, como los incendios, se encuentran vulnerables debido al desconocimiento del entorno. Por ello, es fundamental contemplar estas situaciones. 

\\

Este contraste evidencia la necesidad de soluciones tecnológicas que no solo brinden asistencia en la navegación cotidiana, sino que también se adapten a los diferentes niveles de experiencia, aprendizaje y  necesidades específicas de cada persona. \\

La propuesta de este trabajo terminal plantea una solución que integra dispositivos físicos (lentes inteligentes), visión por computadora, trazado de rayos, sensores de proximidad e inteligencia artificial en un sistema de navegación guiada por voz, especialmente pensado para la movilidad urbana de personas con discapacidad visual. A diferencia de soluciones existentes como aplicaciones móviles o bastones inteligentes, este proyecto propone una interacción manos libres, contextualizada y modular, adaptable a las necesidades individuales del usuario y con potencial de expansión hacia otros sectores vulnerables.

\\

En cuanto a la vinculación con los usuarios potenciales, este proyecto está dirigido principalmente a personas con discapacidad visual, tanto aquellas con ceguera total como parcial, quienes constituyen un sector históricamente desatendido en términos de movilidad autónoma urbana y manejo de riesgos.

Adicionalmente, el sistema podría beneficiar a personas con dificultades temporales de movilidad.\\

La complejidad técnica del proyecto es significativa, ya que involucra el diseño y ensamblado de hardware portátil, el desarrollo de algoritmos de inteligencia artificial, la fusión de datos sensoriales, reconocimiento de objetos tanto estáticos como en movimiento, el trazado de rayos para la reconstrucción de escenarios, la localización y mapeo simultáneo y la creación de una interfaz auditiva accesible.

\\

Todo esto exige la aplicación conjunta de conocimientos adquiridos en asignaturas como visión artificial, aprendizaje automático, procesamiento de 
lenguaje natural, reconocimiento de voz, desarrollo de sistemas embebidos y experiencia de usuario.

\\

Finalmente, el trabajo propuesto es viable dentro del periodo establecido para el desarrollo del TT, considerando que su implementación será gradual y modular siguiendo el método científico, permitiendo validar cada componente de forma independiente. Los recursos requeridos, como sensores de bajo costo, plataformas de desarrollo, y software libre para el entrenamiento de modelos, hacen posible el cumplimiento del proyecto con un presupuesto limitado pero funcional.